% ============================================================
% Related work — Paper A (JNCA): lateral phishing, graph/temporal learning, evaluation rigor
% ============================================================
\section{Related Work}
\label{sec:related}

The review follows the chain of the thesis: lateral phishing and why content is insufficient
(\ref{sec:rw-spear}), the graph and temporal-learning apparatus we borrow (\ref{sec:rw-graph}), and
evaluation rigor (\ref{sec:rw-eval}), closing with the gaps (\ref{sec:rw-gap}).

\subsection{Lateral phishing and BEC: why content is insufficient}
\label{sec:rw-spear}
Targeted attacks---spear phishing, whaling, and Business Email Compromise---differ from mass phishing in their personalization and context.
A recent survey in \emph{Computers \& Security} systematizes the techniques and gaps
of the area \citep{birthriya2025detection}. Content-side classifiers, such as hybrid ensembles over email
content and metadata \citep{bountakas2022helphed}, remain highly effective for the mass tail, but they
inherit the content signal's blindness to clean-content, context-borne deception, and language models now
make personalization cheap \citep{bethany2024evaluating}.
\citet{ho2019detecting} characterized lateral phishing from
compromised accounts across $113$~million emails and fixed an operational low-FPR threshold of $<\!4$~FP/million.
\citet{cidon2019high} targeted payload-free BEC in production, and BEC surveys trace its evolution and cost
\citep{atlam2023business}. What the effective methods share is a departure from content: per-sender behavioral
profiles \citep{gascon2018reading,stringhini2015thataint}, header and stylometry models
\citep{duman2016emailprofiler}, as well as analysis of relay infrastructure \citep{luo2024characterizing} all catch mail
from compromised legitimate accounts whose content reads as normal. Two families of behavioral defense
transfer to a compromised mail account: \emph{per-sender deviation} from the account's own profile, as in COMPA
for social-network accounts \citep{egele2013compa} and in behavioral models of email senders
\citep{stringhini2015thataint}, of which a sudden spike in volume is the simplest form; and \emph{causal-path}
analysis, as in Hopper for lateral movement in enterprise login graphs \citep{ho2021hopper}, which flags a hop to
a rare target without a precursor that explains it. Transferred to email headers, both read the
\emph{sending} side. We show that a stealthy attacker who paces the sends largely removes that signal, and that the receiving side of the
cascade is what remains.

\subsection{Graph and temporal learning for communication graphs}
\label{sec:rw-graph}
The cascade invariant rests on a mature apparatus of graph learning, organized here into the threads that
indicate what we borrow and what is missing for the communication graph of an organization.

On the foundations of inductive learning, modern graph networks learn node representations through
neighborhood aggregation: the \emph{inductive} GraphSAGE \citep{hamilton2017inductive} (generalization to nodes
unseen during training) and the attention-based GAT \citep{velickovic2018graph} establish the paradigm; our \emph{structural} (rather than identity-based) node features are a deliberately
lightweight variant of it, to preserve inductive generalization \citep{xu2020inductive,hao2020inductive}. Link-prediction and scalable-embedding tasks
\citep{zhang2018link,zhang2018scalable} provide heads for classifying $(s\!\to\!v)$ events.

The invariant requires a model of time: the memory-based TGN
\citep{rossi2020temporal}, the attention-based DySAT \citep{sankar2020dysat}, EvolveGCN \citep{pareja2020evolvegcn},
ROLAND \citep{you2022roland}, and streaming and spatio-temporal GNNs
\citep{yu2023streaming,wang2022tsgn,nayebi2025spatiotemporal,behrouz2022anomaly} capture graph evolution.
Temporal graph networks are by now standard instruments for operational network security, from
security-alert prediction on real networks \citep{nayeri2026alert} to real-time anomaly detection on
large-scale provenance graphs \citep{weng2024rtapt}. Our
TGN-lite is a minimal realization of this paradigm by design, allowing the time-shuffle control to isolate the
\emph{mechanism} rather than the capacity of the architecture.

The receiving-side wave has close relatives in the study of diffusion on networks. Outbreak detection asks
where to place sensors so that a spreading process is caught early \citep{leskovec2007outbreak}; source
localization infers the origin of a diffusion from which nodes were reached and when \citep{pinto2012locating,shah2011rumors};
cascade prediction and rumor detection learn from the shape and timing of a spread whether and how far it
will go \citep{cheng2014cascades,vosoughi2018spread,bian2020rumor}. These works model a spreading process on
a known graph, as we do, but they classify the \emph{cascade} or locate its \emph{source}; we classify a
single \emph{event} by the state of its recipient's neighborhood, under a matched design in which every event
is locally legitimate, and we ask specifically whether that signal survives an attacker who paces the
spread.

The closest apparatus comes from fraud detection robust to
\emph{camouflage} and \emph{class imbalance}: CARE-GNN \citep{dou2020enhancing}, GraphConsis \citep{liu2020alleviating},
PC-GNN \citep{liu2021pick}, H2-FDetector \citep{shi2022h2fdetector}, GraphSMOTE \citep{zhao2021graphsmote},
distribution-shift correction \citep{gao2023alleviating,tang2022rethinking}, and the detection of money laundering
and transactional fraud \citep{weber2019antimoney,li2022ttagn,rayana2015collective}. The broader graph-anomaly
literature \citep{ding2019deep,park2020unsupervised,liu2021anomaly,ma2021comprehensive,zhang2021mcgc,li2025umgad}, as well as predicting user dynamics and roles
\citep{kumar2019predicting,wu2020who,li2022pdgnn}, provides techniques, yet operates on graphs of
\emph{transactions/reviews} and has yet to address the \emph{propagation dynamics of a compromise} in a
communication graph. Graph-based phishing detection has so far
targeted mainly \emph{web pages} and \emph{transactions}: graph-theoretic and visual approaches for web pages
\citep{tan2020graph,lin2021phishpedia,abdelnabi2020visualphishnet}, attention over the email graph
\citep{safran2025phishinggnn}, edge-feature GraphSAGE \citep{lo2021egraphsage}, and propagation over URL--IP graphs
\citep{guo2025efficient,li2022pdgnn}. To our knowledge, the receiving-side dynamics of a lateral cascade have
not been operationalized as a detection signal for corporate email.

\subsection{Evaluation rigor: pitfalls and leak control}
\label{sec:rw-eval}
A central caution is that ML detectors in security are easily inflated by leaks, unfair splitting, and headline
AUC in place of operational metrics \citep{arp2022dos}; time-aware evaluation (TESSERACT) shows how temporal
leakage distorts results \citep{pendlebury2019tesseract}. These works motivate our leak-aware protocol and
predictivity probe, which ask not how high the AUC is but whether the result transfers beyond a single corpus and
whether it is an artifact of the generator. We use the same standard to define a method's effectiveness.

\subsection{Gaps and positioning of the work}
\label{sec:rw-gap}
Several gaps emerge from the review. Detection of lateral phishing reads the sending side (volume, paths,
sender profiles), which a patient attacker controls. The propagation dynamics of a compromise on the
organization's communication graph have not been used as a detection signal, and evaluations rarely control for
leaks or test transfer across topologies, so effectiveness tends to be ill-defined. The present work addresses
these gaps: a temporal formalization of the communication graph, the cascade invariant with its causal control,
replication on three real graphs, and a leak-aware methodology that defines effectiveness.
